\scanpage{111}
\mspage{56}

Notons que

\[
\lambda\in Z',\ s\in S\implies\{s,\lambda s\}\notin A_0\tag*{(31)}
\]

C'est clair si $\lambda=\operatorname{id}_S$. Si on avait $\{s,\lambda s\}\in A_0$, soit $t\in S$ tel que $\{s,\lambda s,t\}\in F_0$

\figure{\begin{tikzpicture}[x=1cm,y=1cm]
\draw (0.10,0) -- (1.72,0);
\draw (0,-0.24) -- (0.65,-1.19);
\draw (0.76,-1.20) -- (1.74,-0.14);
\fill (-0.04,0) circle (1.6pt) node[left=3pt] {$s$};
\fill (1.74,0) circle (1.8pt) node[right=4pt] {$\lambda s$};
\node[below] at (0.71,-1.19) {$t$};
\end{tikzpicture}}

et soit $\rho\in G_0^+$ tel que $\deleted{\unsure{g}}\rho(\underbrace{\lambda s}_{\substack{\shortparallel\\\lambda\rho(s)\deleted{\unsure{\lambda}}}})=\lambda s,\ \deleted{g}\rho s=t$, on aura donc $\lambda\rho(s)=\lambda s$ i.e. $\rho(s)=s$ soit $t=s$, absurde.

On suppose maintenant

\[
\begin{aligned}
 (\deleted{\unsure{\mathrm S|8}}\mathrm T_0\,3)\quad&\forall\lambda\in\deleted{\unsure{S}}Z',\ s,t\in S,\ \text{on a}\\
 &\{s,t\}\in A_0\iff\{\lambda s,\lambda^{-1}t\}\in A_0
\end{aligned}\tag*{(32)}
\]

$\bigl[\!\bigl[$\uline{NB} Il suffit de poser \deleted{cet} l'axiome \deleted{pour un}\inserted{$\Longrightarrow$} car il \unsure{est|en} visible, en l'appliquant à $(\lambda^{-1}t,\lambda s)$ au lieu de $(s,t)$, que l'on a

\[
\{\lambda^{-1}t,\lambda s\}\in A_0\implies\deleted{\lambda}\{\underset{\substack{\shortparallel\\\lambda\lambda^{-1}t}}{t},\underset{\substack{\shortparallel\\\lambda^{-1}\lambda s}}{s}\}\in A_0
\]

D'autre part, comme $\lambda,\lambda^{-1}$ centralisent $G_0^+$, il suffit, pour \uline{un} repère origine $\{s_0,t_0,u_0\}$, de vérifier la condition

\[
\{\lambda s_0,\lambda^{-1}t_0\}\in A_0\tag*{(32 bis)}
\]

Si on \deleted{suppose} choisit $\rho\in G_0^+$ tel qu'on ait
% archive sheet number: 110
